\subsection{シンプソン公式の誤差}

単に公式の係数を覚えるだけでは,どれだけ細かく分ければよいかは決められない.ここでは,有限テイラー公式から誤差の定数も導く.座標を移して区間を$[-h,h]$とし,\,$f$は四回連続微分可能とする.\,$-h$での三次テイラー多項式を$P$とすると
\[
 f(x)=P(x)+\frac16\int_{-h}^x(x-t)^3f^{(4)}(t)\dd t.
\]
また,\,$f$の原始関数を$F$,\,$-h$での$F$の四次テイラー多項式を$Q$とする.\,$F'=f$だから$Q'=P$であり,\,$F^{(5)}=f^{(4)}$である.\,$F$の四次の有限テイラー公式を$x=h$で使い,
\[
 F(h)-F(-h)=\int_{-h}^h f(x)\dd x,\qquad
 Q(h)-Q(-h)=\int_{-h}^h P(x)\dd x
\]
を合わせると
\[
 \int_{-h}^hf(x)\dd x
 =\int_{-h}^hP(x)\dd x
  +\frac1{24}\int_{-h}^h(h-t)^4f^{(4)}(t)\dd t
\]
を得る.

一方,最初の式を$x=-h,0,h$で使い,シンプソン公式の重みを掛けて足す.\,$P$に対する誤差は$0$なので,残るのは
\[
 \int_{-h}^hf(x)\dd x
 -\frac h3\{f(-h)+4f(0)+f(h)\}
 =\int_{-h}^hK(t)f^{(4)}(t)\dd t
\]
である.ただし,係数をまとめた$K$は
\[
 K(t)=
 \begin{cases}
 -\frac{(h+t)^3(h-3t)}{72}&(-h\leq t\leq0),\\[4pt]
 -\frac{(h-t)^3(h+3t)}{72}&(0\leq t\leq h).
 \end{cases}
\]
例えば$t\geq0$では,中点$f(0)$の剰余項に$t$は現れないので,係数は
\[
 \frac{(h-t)^4}{24}-\frac h{18}(h-t)^3
 =-\frac{(h-t)^3(h+3t)}{72}
\]
となる.\,$t\leq0$では,これに中点からの項$-\frac{2h}{9}(-t)^3$を加え,整理すると上の式になる.

$K$は偶関数であり,全区間で$K\leq0$である.その積分の絶対値は
\[
\begin{aligned}
 \int_{-h}^h|K(t)|\dd t
 &=\frac1{36}\int_0^h(h-t)^3(h+3t)\dd t\\
 &=\frac{h^5}{36}\int_0^1(1-u)^3(1+3u)\dd u
 =\frac{h^5}{90}.
\end{aligned}
\]
したがって$|f^{(4)}|\leq M_4$なら,一組の区間について
\[
 \left|\int_{m-h}^{m+h}f(x)\dd x
 -\frac h3\{f(m-h)+4f(m)+f(m+h)\}\right|
 \leq\frac{M_4h^5}{90}
\]
となる.さらに$f^{(4)}\geq0$なら,積分の値から近似値を引いた誤差は$0$以下である.

\begin{theorem}{合成シンプソン公式と誤差}
$N$を正の偶数とし,\,$h=\frac{b-a}{N}$,\,$x_j=a+jh$とする.\,$f$が$[a,b]$で四回連続微分可能で,\,$|f^{(4)}(x)|\leq M_4$なら
\[
 S_N=\frac h3\left\{f(a)+f(b)
 +4\sum_{j=1}^{\frac N2}f(x_{2j-1})
 +2\sum_{j=1}^{\frac N2-1}f(x_{2j})\right\}
\]
に対して
\[
 \left|\int_a^bf(x)\dd x-S_N\right|
 \leq\frac{M_4(b-a)h^4}{180}
 =\frac{M_4(b-a)^5}{180N^4}
\]
が成り立つ.
\end{theorem}

小区間を二つずつ組にすると組数は$\frac N2$である.各組の誤差上限$\frac{M_4h^5}{90}$を足すことで,この評価が得られる.組どうしの境目は左右の組で一回ずつ使われるため係数が$2$,各組の中点は係数が$4$になる.分割数を二倍にすると誤差上限は十六分の一になるが,これは四階微分が区間全体で有界という条件の下での結論である.端点で微分が発散する関数に,そのままこの上限を使うことはできない.

\begin{example}{近似公式の比較}
$\int_0^1\frac1{1+x}\dd x$について,四等分した台形公式とシンプソン公式を計算し,誤差を評価せよ.
\end{example}

分点は$0,\frac14,\frac12,\frac34,1$で,関数値は順に$1,\frac45,\frac23,\frac47,\frac12$である.したがって
\[
 T_4=\frac14\left(\frac34+\frac45+\frac23+\frac47\right)
 =\frac{1171}{1680},
\]
\[
 S_4=\frac1{12}\left\{\frac32+4\left(\frac45+\frac47\right)+\frac43\right\}
 =\frac{1747}{2520}.
\]
二階微分は$\frac2{(1+x)^3}$,四階微分は$\frac{24}{(1+x)^5}$で,いずれも正である.そのため,どちらの近似値も真の値以上になる.上限$M_2=2$,\,$M_4=24$を使えば
\[
 T_4-\frac1{96}\leq\int_0^1\frac1{1+x}\dd x\leq T_4,
\]
\[
 S_4-\frac1{1920}\leq\int_0^1\frac1{1+x}\dd x\leq S_4.
\]
この積分は$\log2$に等しいが,その小数値を知らなくても近似の保証は作れる.同じ五つの関数値を使っていても,曲線をどの次数の式で置き換えるかによって誤差の評価が変わる.



